% ── Section V-A: Single-Flow Handover Performance ────────────────────────────
% Placed before the fairness section in the Evaluation.

\subsection{Single-Flow Handover Performance}
\label{sec:exp1}

\subsubsection{Setup}

We drive each of the five baselines through all six pairwise orbit
transitions (LEO$\leftrightarrow$MEO, LEO$\leftrightarrow$GEO,
MEO$\leftrightarrow$GEO) at three handover times
($T_\mathrm{HO} \in \{30, 60, 120\}$\,s) and seven lead times
($\ell \in \{0, 2, 5, 10, 20, 30, 60\}$\,s).
Each condition is a zero-loss, single-flow simulation; parameters for each
orbit are given in Table~\ref{tab:orbits}.
The primary metric is \textbf{T90}: the elapsed time from
$T_\mathrm{HO}$ until the delivered throughput first reaches
90\,\% of the new orbit's link capacity, measured as a one-second rolling
average.
A trial is declared \emph{not converged} (N/C) if T90 exceeds the
60-second post-handover window.
We additionally record post-handover steady-state throughput
(mean over $t \in (T_\mathrm{HO}+20, T_\mathrm{HO}+60]$\,s),
goodput (total bytes delivered in the 90-second run),
and peak on-path queue depth.

\subsubsection{LEO$\to$GEO: the critical transition}

A downward BDP transition --- lower link rate, higher latency, larger BDP
--- is the most demanding case.
Table~\ref{tab:exp1_leo_geo} and Figure~\ref{fig:t90_lead} show results
for the LEO$\to$GEO transition at $T_\mathrm{HO} = 30$\,s with $\ell = 0$.

\textbf{B1 (vanilla BBRv3)} fails entirely: the on-path queue inflates to
33\,MB (the full GEO buffer), steady-state throughput is effectively zero,
and the 90\% threshold is never reached within the 60-second window.
The root cause is that BBRv3 retains its LEO-calibrated bandwidth estimate
after handover.
Because BBRv3's two-cycle \texttt{MaxBwFilter} clears only after roughly
$8 \times \mathrm{RTT}_\mathrm{GEO} \approx 5$\,s~\cite{cardwell2017},
the sender continues pacing at $\approx$50\,Mbps into a 10\,Mbps pipe,
sustaining a queue that grows without bound at the simulation's buffer limit.

\textbf{B3 (cwnd-freeze)} exhibits the same failure mode.
Freezing the congestion window at the pre-handover LEO value leaves an
equally oversized window in place; the buffer fills to 33\,MB
identically to B1.
Advance warning is useless: even at $\ell = 30$\,s, B3 does not converge.

\textbf{B4 (send-pause / resume)} converges exactly once across the full
lead-time sweep, at $\ell = 5$\,s (T90 = 5.6\,s), and fails at every
other lead time.
The narrow success window reflects the tight timing between the pause
duration and the GEO queue-drain time: too short a pause leaves residual
in-flight data that re-inflates the queue; too long a pause wastes capacity
and triggers a slow-start on resume.

\textbf{BBR-SAT} converges at every tested lead time from $\ell = 0$ to
$\ell = 20$\,s, with T90 = 2.5\,s throughout.
On receiving the CONFIRMED signal at $T_\mathrm{HO}$, the algorithm
loads the seeded GEO orbit entry (\texttt{bw\_hi} = 10\,Mbps, \texttt{min\_rtt}
= 580\,ms), enters a brief DRAIN phase to clear residual in-flight data,
then transitions to REFILL and converges to a steady-state of 9.2\,Mbps
(92\,\% of link capacity) with a peak queue of only 910\,KB --- a
36$\times$ reduction compared to B1/B3.
Goodput over the 90-second run is 56.9\,MB, largely independent of lead time.

At $\ell = 30$\,s, T90 degrades sharply to 47.6\,s.
The PREDICTED signal at $t = 0$\,s (i.e.\ $T_\mathrm{HO} - 30$\,s)
initiates a proactive queue drain via ProbeBW\_DOWN; however, because
the actual handover does not occur for 30\,s, BBR-SAT re-enters
ProbeBW\_UP and rebuilds a LEO-calibrated queue before the CONFIRMED
signal fires.
This establishes a practical upper bound: lead times up to approximately
20\,s are beneficial; beyond that, pre-draining fires too early to hold.

\textbf{CUBIC} converges in 1.5\,s at all lead times, achieving a
steady-state of 10.0\,Mbps (100\,\% utilization) and 66.4\,MB goodput.
The rapid convergence is loss-driven: the queue overflow immediately after
handover triggers multiplicative window reduction, and CUBIC settles within
two to three 580\,ms GEO RTTs.
Although CUBIC achieves slightly higher single-flow throughput than
BBR-SAT (at the cost of a 1.3\,MB peak queue versus 910\,KB), its
behaviour under competition is discussed in Section~\ref{sec:fairness}.

\begin{table}[t]
\centering
\caption{LEO$\to$GEO, $T_\mathrm{HO}=30$\,s, $\ell=0$, no loss}
\label{tab:exp1_leo_geo}
\small
\begin{tabular}{lrrrr}
\toprule
Baseline & T90 & SS bw & Goodput & Peak queue \\
\midrule
B1 Vanilla BBRv3  & N/C   & 0.0\,Mbps & 5.4\,MB  & 33\,120\,KB \\
B3 cwnd-freeze    & N/C   & 0.0\,Mbps & 5.4\,MB  & 33\,120\,KB \\
B4 pause/resume   & N/C   & 0.0\,Mbps & 3.8\,MB  & 34\,238\,KB \\
BBR-SAT (ours)    & \textbf{2.5\,s} & \textbf{9.2\,Mbps} & \textbf{56.9\,MB} & \textbf{910\,KB} \\
B5 CUBIC          & 1.5\,s & 10.0\,Mbps & 66.4\,MB & 1\,328\,KB \\
\bottomrule
\end{tabular}
\end{table}

\subsubsection{Lead-time sensitivity}

Figure~\ref{fig:t90_lead} plots T90 as a function of lead time for the
LEO$\to$GEO transition.
BBR-SAT maintains a flat T90 of 2.5\,s across $\ell \in [0, 20]$\,s,
confirming that the PREDICTED signal is effectively used but that even
zero lead time is sufficient for convergence given the seeded orbit table.
The practical interpretation is that any advance notice in the range
2--20\,s is equally effective; the value of the PREDICTED signal lies
primarily in \emph{proactive queue drain} that prevents mid-flight data
from inflating the post-handover queue, and this drain completes within
one to two GEO RTTs.

\subsubsection{BW-increasing and moderate transitions}

For all bandwidth-increasing transitions (MEO$\to$LEO, GEO$\to$LEO,
GEO$\to$MEO) and the moderate MEO$\to$GEO and LEO$\to$MEO transitions,
every baseline converges at $\ell = 0$ with T90 $\leq 3.5$\,s
(Table~\ref{tab:exp1_all_transitions}).
Bandwidth increases are inherently self-correcting: the new link can absorb
the pre-handover cwnd without queue buildup, and BBR's bandwidth probing
detects the higher capacity within one probe cycle.
BBR-SAT's T90 at $\ell = 0$ is 0--1\,s higher than B1/B3 on these
transitions because the CONFIRMED signal briefly forces a REFILL
phase before resuming ProbeBW; with $\ell \geq 2$\,s this overhead
disappears as the PREDICTED signal pre-positions the orbit context.

\begin{table}[t]
\centering
\caption{T90 at $\ell = 0$, $T_\mathrm{HO} = 30$\,s, no loss (all transitions)}
\label{tab:exp1_all_transitions}
\small
\setlength{\tabcolsep}{4pt}
\begin{tabular}{lrrrrr}
\toprule
Transition & B1 & B3 & B4 & BBR-SAT & CUBIC \\
\midrule
LEO$\to$MEO  & 0.5\,s & 0.5\,s & 6.5\,s & 8.5\,s  & 0.5\,s \\
LEO$\to$GEO  & N/C    & N/C    & N/C    & \textbf{2.5\,s} & 1.5\,s \\
MEO$\to$LEO  & 2.5\,s & 2.5\,s & 0.5\,s & 3.5\,s  & 0.5\,s \\
GEO$\to$LEO  & 1.5\,s & 1.5\,s & 2.5\,s & 2.5\,s  & 1.5\,s \\
MEO$\to$GEO  & 1.5\,s & 1.5\,s & 1.5\,s & 2.5\,s  & 1.5\,s \\
GEO$\to$MEO  & 1.5\,s & 1.5\,s & 2.5\,s & 2.5\,s  & 1.5\,s \\
\bottomrule
\end{tabular}
\end{table}

\subsubsection{Handover-time sensitivity}

Repeating the LEO$\to$GEO experiment at $T_\mathrm{HO} = 60$\,s reveals
moderate T90 inflation: BBR-SAT achieves 2.5\,s at $\ell \in \{0, 20\}$\,s
but 7.5--9.5\,s at intermediate lead times, likely due to the orbit table's
minimum-RTT entry drifting over the longer pre-handover interval as
BBR natural probing samples slightly elevated RTTs.
At $T_\mathrm{HO} = 120$\,s all BBR-SAT conditions fail to converge,
suggesting that the orbit table's 10-second \texttt{min\_rtt} expiry window
allows the seeded GEO RTT entry to become stale before the handover fires.
Extending the orbit table's RTT confidence window for long-duration pre-handover
periods is identified as future work.
